% !TEX root = ../nc-cubic-n-fold.tex
% Tate--Koszul reconstruction of origin/rotation-preserving enhanced kernels.

\subsection{原点と内部 shift を保つ核の Tate--Koszul 復元}

以下では $B$ を標数零の体 $k$ 上の連結局所有限な左右 noetherian
quadratic Koszul AS-regular 代数とし、$\Omega\in Z(B)_3$ を
非零の正則元とする。$A=B/(\Omega)$ と置く。
$\kappa$ は $k=A/A_{\ge1}$ の graded singularity category における像、
$\rho=(1)$ は内部 shift である。
ここで取り扱うのは enhancement を保つ perfect 核であり、
三角関手の自然同型だけを dg 自然同型と置き換えない。
また、$\rho$ と可換であるという性質と、指定された
$\rho^3\simeq[2]$ の標識を保つことは区別する。

\begin{lemma}[中央三次商の通常 Ext と Tate Ext]
\label{lem:nfm-cubic-tate-cut}
双次数を
\[
 E^{p,q}=\operatorname{Ext}^p_{\operatorname{Gr}A}(k,k(q)),
 \qquad
 S^{p,q}=\operatorname{Hom}_{D_{\mathrm{sg}}^{\mathbb Z}(A)}
                 (\kappa,\kappa(q)[p])
\]
とする。指定された $\Omega$ に由来する中央 class
$t\in E^{2,-3}$ に関し、積を含む同型
\[
 E\simeq B^!\otimes_k k[t],\qquad
 S\simeq B^!\otimes_k k[t,t^{-1}]
\]
がある。$B^!_i$ の双次数は $(i,-i)$ である。
さらに実際の stabilization が誘導する Adams graded dg algebra の射
\[
 \mathcal E=R\underline{\operatorname{End}}_A(k)
       \longrightarrow
 \mathcal S=R\underline{\operatorname{End}}_{D_{\mathrm{sg}}(A)}(\kappa)
\]
は、$t$ に関する derived localization を表示する。
\end{lemma}

\begin{proof}
中央正則元の長さ一の分解から、$A\otimes_B^{\mathbf L}A$ の cohomology は
$H^0=A$、$H^{-1}=A(-3)$ だけである。$A$-双加群として標準 truncation を取ると、
本稿の基礎節の命題\ref{lp:prop:sing-period}
と同じ自己交叉三角
\[
 A(-3)[1]\longrightarrow A\otimes_B^{\mathbf L}A
 \longrightarrow A\xrightarrow{\eta}A(-3)[2]
\]
を得る。これは実際の $A$-双加群の三角である。
任意の有界有限生成 $A$-複体 $M$ に評価すると
\begin{equation}\label{eq:nfm-tate-self-intersection}
M(-3)[1]\longrightarrow M\otimes_B^{\mathbf L}A
 \longrightarrow M\xrightarrow{\eta_M}M(-3)[2].
\end{equation}
$B$ の有限大域次元により中項は $A$-perfect である。
$t=\eta_k$ は双加群射に由来するので、全 graded Ext 上で中央に作用する。
内部 shift の規約からその双次数は $(2,-3)$ である。

$M=k$ に対し $R\operatorname{Hom}_A(-,k(q))$ を適用する。
導来随伴によって中項の Ext は $\operatorname{Ext}_B(k,k(q))$ であり、
長完全列は
\begin{equation}\label{eq:nfm-tate-change-rings}
\cdots\longrightarrow E^{p-2,q+3}\xrightarrow{t}E^{p,q}
 \longrightarrow\operatorname{Ext}_B^{p,q}(k,k)
 \longrightarrow E^{p-1,q+3}\longrightarrow\cdots
\end{equation}
となる。連結次数一生成代数の minimal 分解より
$E^{r,s}=0$ for $s>-r$ である。
一方 $B$ は Koszul なので $\operatorname{Ext}_B^{p,q}=0$ unless $q=-p$。
従って右の connecting map の非零を考える必要があるのは
$q=-p$ の場合だけだが、その終域は
$E^{p-1,3-p}=0$ である。左側の connecting も同じ理由で零である。
よって各双次数で
\begin{equation}\label{eq:nfm-tate-short-exact}
0\longrightarrow E^{p-2,q+3}\xrightarrow{t}E^{p,q}
 \longrightarrow\operatorname{Ext}_B^{p,q}(k,k)\longrightarrow0.
\end{equation}
特に $q=-p$ の対角では最初の項も零なので restriction は同型である。
この restriction は導来 restriction functor が誘導する Yoneda algebra の射である。
対角成分は積で閉じているので、その逆が $B^!$ を $E$ の部分代数として与える。
$t$ の中央性と \eqref{eq:nfm-tate-short-exact} の反復から、通常 Ext の表示を得る。

次に \eqref{eq:nfm-tate-self-intersection} の $\eta$ を反転する。
任意の $A$-perfect 複体 $P$ に対し、十分大きい $r$ では
$\operatorname{Hom}(P,P(-3r)[2r])=0$ だから $\eta_P^r=0$ である。
逆に $\eta_M^r=0$ なら、その cone は \eqref{eq:nfm-tate-self-intersection} と octahedral triangle により perfect であり、
零写像の cone の直和因子である $M$ も perfect となる。
従って $\eta$-torsion 対象はちょうど perfect 対象である。
中央自然変換の Verdier localization は
\begin{equation}\label{eq:nfm-tate-hom-localization}
\operatorname{Hom}_{D_{\mathrm{sg}}(A)}(M,N(q)[p])
 \simeq\underset{r}{\operatorname{colim}}\,
 \operatorname{Hom}_{D^b(\operatorname{gr}A)}
       (M,N(q-3r)[p+2r])
\end{equation}
を与える。具体的には roof の左矢印の cone が perfect であるとき、
その cone を消す $\eta$ の冪を用いて roof を右辺の射で表せる。
同様に roof の零性も、さらに大きい冪の後で検出される。
これで通常 Ext から Tate Ext への写像は $t$-局所化である。

dg algebra の主張には、Braun--Chuang--Lazarev の derived localization
\cite[Theorems~3.18 and~5.3]{NFMBCLLocalization} を適用する。
$t$ の像は $H(\mathcal S)$ で可逆だから
$L_t\mathcal E\to\mathcal S$ が存在し、同文献 Theorem~5.3 と \eqref{eq:nfm-tate-hom-localization} により
cohomology 上で同型となる。ここでは Adams 次数ごとに斉次生成元を持つ
semi-free/cofibrant resolution を取り、局所化の可逆元とその homotopy を
$t$ の内部次数に合わせて付加する。従って同じ derived localization の構成と普遍性が
Adams graded 複体でも成立し、同定は内部次数を保つ。
これは cohomology algebra だけの同定ではない。
\end{proof}

\begin{lemma}[Tate 模型内の通常 Ext の回復]
\label{lem:nfm-ordinary-ext-intrinsic-cut}
$\mathcal S$ の strictly unital な双次数付き minimal $A_\infty$ 模型を $S$ とする。
\[
 j=-p-q,\qquad
 S_{\ge0}=\bigoplus_{-p-q\ge0}S^{p,q}
\]
と置く。この部分空間は augmented $A_\infty$ 部分代数であり、
$\mathcal E$ の minimal 模型と $A_\infty$ 同型である。
任意の内部次数を保つ $A_\infty$ 自己同型は $S_{\ge0}$ を保ち、
その homotopy もこの部分に制限できる。
\end{lemma}

\begin{proof}
前補題の表示では、$B^!_i t^j$ の双次数は $(i+2j,-i-3j)$ なので、
上の $j$ は実際の $t$ 指数に一致する。
minimal 構造の $m_r$、射の $f_r$、reduced homotopy の $h_r$ は
それぞれ cohomological degree $2-r,1-r,-r$、internal degree $0$ である。
従って、入力の $t$ 指数を $j_1,\ldots,j_r$ とすると出力の指数は順に
\begin{equation}\label{eq:nfm-tate-degree-cutoff}
\sum j_i+r-2,\qquad \sum j_i+r-1,\qquad \sum j_i+r.
\end{equation}
$m_1=0$ と $m_0=0$ を考慮すれば、全て $j_i\ge0$ を保つ。
非単位元の内部次数は負なので augmentation も保たれる。

stabilization の dg 射を双次数付き minimal 模型へ移す。
その $A_\infty$ Taylor 成分も \eqref{eq:nfm-tate-degree-cutoff} を満たすので $S_{\ge0}$ に因子化する。
前補題より一次成分は同型である。
minimal $A_\infty$ 代数間の一次成分が同型な射は、bar coalgebra の長さに関する
再帰により $A_\infty$ 逆を持つ。これが所要の比較である。
homotopy の制限についても \eqref{eq:nfm-tate-degree-cutoff} の第三式を使えばよい。
\end{proof}

\begin{proposition}[通常 Ext への制限の忠実性]
\label{prop:nfm-tate-automorphism-faithful}
内部次数を保つ strictly unital $A_\infty$ 自己同型の homotopy class から、
次数付き代数 $A$ の自己同型への反変単射
\begin{equation}\label{eq:nfm-tate-automorphism-restriction}
\pi_0\operatorname{Aut}_{A_\infty,\mathrm{Adams}}(S)
 \lhook\joinrel\longrightarrow\operatorname{Aut}_{\mathrm{gr}}(A)
\end{equation}
がある。反変という向きを逆元によって直せば群準同型となる。
特に、右辺で恒等を誘導する自己同型は Tate 模型上でも homotopic to identity である。
\end{proposition}

\begin{proof}
前補題によって通常 Ext へ制限する。
Lu--Palmieri--Wu--Zhang\cite[Theorem~2.4 and Corollary~2.5(b)]{NFMLPWZKoszul}
の Adams Koszul 双対は、augmented $A_\infty$ homotopy category の Hom を反変に全単射で移す。
ここで必要な局所有限性は全て Adams connected 性から成り立つ。
実際、固定した負の内部次数に現れる Ext は有限次元で、
cohomological degree も有限区間に入る。bar 構成でもその内部次数での長さは有限である。
元の $A$ は cohomological degree $0$ に集中しているので、
$A\to A$ の $A_\infty$ 射は degree 理由で strict な代数射であり、
その間の reduced homotopy も零である。
従って通常 Ext の自己同型 class は $\operatorname{Aut}_{\mathrm{gr}}(A)$ へ反変に同型で移る。
ここで完備化した代数に置き換えてはいない。

制限後に二つの Tate 自己同型が homotopic になったとする。
通常 Ext から $S$ への二つの合成は homotopic である。
前々補題の derived localization と
Braun--Chuang--Lazarev\cite[Proposition~3.17 and Theorem~3.18]{NFMBCLLocalization} により、
局所化射による derived mapping space の restriction の fibre は空または可縮である。
従ってこの二つの $S$ 上の射自体が homotopic である。
これが \eqref{eq:nfm-tate-automorphism-restriction} の単射性である。
一般の pro-unipotent 群の homotopy 商についての閉性や一意可除性を、
この証明で別に仮定する必要はない。
\end{proof}

\begin{remark}[三次関係と高次 gauge]
\label{rem:nfm-cubic-relations-modulo}
$V=A_1=B_1$、$R\subset V^{\otimes2}$ を二次関係とする。
任意の次数付き $A$ 自己同型は $V$ 上の線型作用によって決まり、
$R$ を保つので $B=T(V)/(R)$ の自己同型を誘導する。
三次関係空間は
\[
 \ker(V^{\otimes3}\to A_3)
   =RV+VR+k\widetilde\Omega
\]
である。従ってこの自己同型は $B_3$ の直線 $k\Omega$ を保ち、逆も成り立つ。
\begin{equation}\label{eq:nfm-tate-relation-automorphisms}
\operatorname{Aut}_{\mathrm{gr}}(A)
  =\operatorname{Aut}_{\mathrm{gr}}(B,k\Omega).
\end{equation}
Lu--Palmieri--Wu--Zhang\cite[Corollary~B]{NFMLPWZExt} の高次積の記述では、
$m_2$ と $m_3$ がこの二次・三次関係を回復する。
ただし $m_3$ の三次テンソルそのものを gauge 不変と主張してはいけない。
$f_2(\xi,\xi)$ の出力 $(1,-2)$ は零だが、
$f_2(m_2(\xi,\xi),\xi)$ の出力 $(2,-3)$ は零とは限らず、
$m_3$ は $RV+VR$ を法としてのみ $\Omega$ を定める。
\eqref{eq:nfm-tate-automorphism-restriction}--\eqref{eq:nfm-tate-relation-automorphisms} はこの gauge の曖昧さを quotient の代数 $A$ によって処理している。
\end{remark}

\begin{lemma}[内部 shift の比較同型と graded orbit algebra]
\label{lem:nfm-kernel-orbit-algebra}
$\mathcal D=\operatorname{thick}\{\kappa(i):i\in\mathbf Z\}$ とし、
$HH^0(\mathcal D)=k$ を仮定する。
実際の enhanced 自己同値 $Q$ が
$Q\kappa\simeq\kappa$ および $Q\rho\simeq\rho Q$ を満たすとき、
上の $S$ の Adams graded $A_\infty$ 自己同型を誘導する。
$Q$ の同型類だけからは、その自己同型 class は内部次数の character
\begin{equation}\label{eq:nfm-tate-character}
a\in S^{p,q}\longmapsto c^q a\qquad(c\in k^\times)
\end{equation}
だけの曖昧さを持つ。この曖昧さを除いた作用が恒等なら、$Q$ 自体も恒等核の class である。
\end{lemma}

\begin{proof}
まず $\theta:Q\kappa\xrightarrow{\sim}\kappa$ と
実際の quasi-functor の同型 $\eta:Q\rho\xrightarrow{\sim}\rho Q$ を一つ選ぶ。
\emph{この段階では指定三次周期の図式を課さない。}
一つの $\eta$ の正負の反復が $\mathbf Z$-作用との比較を与える。
$\mathbf Z$ は一つの可逆生成元からなるので、別の relation の coherence は課されない。
cofibrant bimodule を使えば、この操作は Keller の dg orbit 構成と
その functoriality\cite[Sections~9.3--9.4]{NFMKellerOrbit} により
実際の enhanced 関手となる。
この構成で orbit の次数 $q$ は保持する。
cofibrant replacement と bar resolution も $q$ 次数付きで取る。

原点でこの関手の endomorphism を取ると
\[
 \bigoplus_q R\operatorname{Hom}_{\mathcal D}(\kappa,\rho^q\kappa)
\]
の Adams graded $A_\infty$ 自己同型を得る。
より具体的には、$\theta$ と $\eta$ を反復して
$\theta_i:Q\rho^i\kappa\simeq\rho^i\kappa$ を作り、
全対象 $\rho^i\kappa$ の Hom の作用を移す。
その相対次数 $j-i$ が上の $q$ である。
Keller の構成は pair $(Q,\eta)$ の合成と両立するので、
この操作は自己同値の合成も保つ。

前補題の cohomology 計算から
$\operatorname{End}^0(\kappa)=S^{0,0}=k$ である。
従って $\theta$ の変更は scalar の共役だけで、作用を変えない。
また二つの $\eta$ の比は $Q\rho$ の scalar 自己同型である。
これは $HH^0(\mathcal D)=k$ と可逆性による。
$\eta$ を $c\eta$ にすると、$q$ 回の反復には $c^q$ が現れ、\eqref{eq:nfm-tate-character} だけが変わる。
選んだ同型の chain lift や homotopy の変更は derived な同じ射を表すので、
得られた $A_\infty$ map の homotopy class を変えない。

最後に忠実性を確認する。この際、\emph{Adams 次数を忘れた orbit 圏だけ}は使わない。
双次数付き $S$ から、対象集合 $\mathbf Z$、Hom$(i,j)=S_{j-i}$ の
$A_\infty$ 圏を作る。この圏は、$\mathcal D$ の全対象
$\{\rho^i\kappa\}$ からなる充満部分圏の minimal 模型である。
Adams 次数を保つ homotopy はこの圏の各 Hom に持ち上がる。
従って $S$ 上の作用が homotopic to identity なら、$Q$ のこの充満部分圏への制限も
恒等関手に同型である。
この部分圏は $\mathcal D$ を生成するので、pretriangulated hull と idempotent completion
への Morita の普遍性によって $Q\simeq\mathrm{id}_{\mathcal D}$ である。
ここで得るのは enhanced 関手、従って実際の可逆 bimodule/kernel の同型である。
\end{proof}

\begin{theorem}[原点・rotation 安定化群の表示自己同型への忠実性]
\label{thm:nfm-origin-rotation-faithful}
$\mathcal D=\operatorname{thick}\{\kappa(i):i\in\mathbf Z\}$ とし、
$HH^0(\mathcal D)=k$ を仮定する。原点の class と内部 shift を保つ
$\mathcal D$ の enhanced 自己同値の群には単射
\begin{equation}\label{eq:nfm-origin-rotation-faithful}
\{[Q]:Q\kappa\simeq\kappa, Q\rho\simeq\rho Q\}
 \lhook\joinrel\longrightarrow
 \operatorname{Aut}_{\mathrm{gr}}(B,k\Omega)/\mathbf G_m
\end{equation}
がある。右辺の $\mathbf G_m$ は次数スカラーである。
特に、この表示自己同型群が自明なら左辺も自明である。
\end{theorem}

\begin{proof}
前補題によって $S$ 上の作用を取り、命題
\ref{prop:nfm-tate-automorphism-faithful} と \eqref{eq:nfm-tate-relation-automorphisms} を適用する。
反変の向きは inverse を取って群準同型に直す。
\eqref{eq:nfm-tate-character} は $A_1$ 上で scalar に移るので、右辺の次数スカラーで商を取れば
$\eta$ の選択によらない写像になる。
核に入る $Q$ では、$\eta$ を scalar 倍して $A$ 上の作用を恒等にできる。
命題\ref{prop:nfm-tate-automorphism-faithful} により Tate 上の作用は homotopic to identity であり、
前補題の最後の忠実性から $Q$ は恒等核の class である。
有限群の coherent action の存在や平均化を使っていないので、
有限 kernel の捩れを別に除く必要はない。
\end{proof}

\begin{lemma}[一般一ブロック対の表示安定化子]
\label{lem:nfm-generic-presentation-stabilizer-trivial}
良い一ブロック表示には非空不変開があり、その任意の幾何点で
\[
 \operatorname{Aut}_{\mathrm{gr}}(B,k\Omega)/\mathbf G_m=1
\]
である。ここで対象は周囲代数 $B$ 単独でなく中心直線を含む対であり、
$\mathbf G_m$ は次数スカラーである。
\end{lemma}

\begin{proof}
命題\ref{prop:nfm-graded-presentation}の全次数付き同型の表示と
定理\ref{thm:poisson-good-open-isomorphism}により、この群は対応する
Poisson 対 $(X,\pi)$ の自己同型群と一致する。後者は、
$X=V(H_\Omega)$ の射影自己同型群の部分群である。
ここで $H_\Omega$ は trace--Hessian の逆構成で得られる滑らかな三次であり、
残余三次との同一視を仮定していない。

系\ref{cor:poisson-generic-open-immersion}により、この良い対の族は
滑らかな三次のモジュライを支配する。
$n\ge3$ では一般の滑らかな三次 $n$-fold の射影自己同型群は自明であり、
その条件は非空不変開を定める
\cite[Theorem~1.5 and the preceding paragraph]{NFMPoonenAutomorphisms}。
その開を引き戻せば所要の非空開となる。
次数付き自己同型の射影化の核は、次数一生成性により次数スカラーだけである。
\end{proof}

\begin{corollary}[全次元の正準な随伴復元点型 class]
\label{cor:nfm-generic-adjoint-class-unique}
固定入力の良い標準 split 一ブロック族の一般係数体 $K=k_n$ を考える。
全 $n\ge3$ で、定理\ref{thm:nfm-adjoint-point-finite-class-torsor}の群と torsor は
\[
 H_p=1,\qquad S_p=\operatorname{Spec}K
\]
を満たす。従って随伴復元条件を満たす点型 full 比較核の幾何学的 class は
一つに定まり、$K$ 上で正準的に選択される。
その実現 gerbe は
\[
 \mathscr G_p\simeq\mathscr P_p,
 \qquad [\mathscr G_p]=\delta_C(p),
 \qquad \operatorname{per}\delta_C(p)=\operatorname{ind}\delta_C(p)=3
\]
であり、この型の実際の核の最小有限実現次数は三である。
ここでは rotation と可換であることを要求するが、比較同型と指定三次周期の標識は
実現 stack の追加データに含めていない。
\end{corollary}

\begin{proof}
適用する代数的仮定と圏の範囲を先に確認する。
本稿の基礎節の
\texttt{field\_model.tex}, \texttt{gcd:prop:automatic-as} は、
次数別有限自由な半古典模型の一般 ambient が左右 noetherian な二次 Koszul
AS-regular 代数であり、指定中心三次元が正則であることを与える。
その証明の入力は \texttt{local\_model.tex} の
\texttt{lp:lem:noetherian}, \texttt{lp:prop:koszul},
\texttt{lp:prop:gorenstein} である。
良い一ブロック表示の一般体はこの支配的な半古典族へ埋め込まれ、
これらの次数付き有限自由分解と正則性は忠実な体拡大から降下する。
従って本節の $B,A=B/(\Omega)$ に関する仮定を満たす。

同入力 \texttt{window.tex}, \texttt{tw:lem:stabilization} は
実際の dg stabilization と singularity category の Morita 比較を与える。
\texttt{tw:thm:window}, \texttt{tw:eq:main} により、source の線束生成残余圏は
\[
 \operatorname{thick}\{\kappa,\kappa(1),\kappa(2)\}
 =\operatorname{thick}\{\kappa(i):i\in\mathbf Z\}
\]
と同値であり、rotation は内部 shift へ移る。

coherent 飽和には \texttt{coherent\_saturation.tex} の
\texttt{crs:thm:coherent} を用いる。
一般点の比較は \texttt{main\_result.tex} の
\texttt{gdr:thm:generic} による。
これらにより、上の圏は一般点で全 coherent 残余圏に一致する。
従って生成対象の部分圏だけについての結論を全圏へ無条件に拡大しているのではない。

随伴復元点型の幾何学的比較核 $P_0$ を一つ取る。
定理\ref{thm:nfm-adjoint-point-finite-class-torsor}の非空性と
$\Phi_{P_0}(\kappa[-1])\simeq E_p$、rotation 適合性から、
$P_0$ による共役は $H_p$ を原点と内部 shift を保つ source 自己同値群に移す。
命題\ref{prop:nf-hh}により $HH^0=k$ である。
定理\ref{thm:nfm-origin-rotation-faithful}と前補題により、この群は自明である。
よって $H_p=1$、非空なその torsor は $\operatorname{Spec}K$ となる。
$P_0$ の選択は証明にのみ使い、得られる唯一の class はその選択によらない。

実現 gerbe の式は定理\ref{thm:nfm-adjoint-point-finite-class-torsor}の積表示による。
標準 split 一般係数体での $\delta_C(p)$ の exact index 三は、
系\ref{cor:nfm-generic-residual-picard-index-three}により
点側の非消滅と半古典平面同型から従う。
従って系\ref{cor:nfm-adjoint-point-realization-degree}により最小次数は三である。
\end{proof}

\begin{remark}[復元の忠実性と残る範囲]
上の証明は、原点の対象像を固定しただけで endomorphism の全作用が恒等になると
仮定していない。Tate 模型から通常 Ext を双次数で切り出し、
Koszul 双対と derived localization の二つの忠実性を通してその作用を検出した。
特に $m_3$ の gauge を無視した三次テンソルの復元や、
homotopy 恒等部分群の pro-unipotence は必要ない。

また、内部 shift との比較同型の scalar は Adams character として追跡し、
次数スカラーで割ってから核の class に戻った。
指定された三次周期の coherence をここで新たに選択したのではなく、
追加の $\mu_3$ 位相の消滅はこの一意性からは従わない。
随伴復元条件を持たない任意の full 核の全群分類、
他の intrinsic 成分の分類、coherent な中心線型性も別の問題である。
\end{remark}

\subsection{随伴復元対象の低次積による周期位相の有限評価}

本節では標準 split 一ブロックの一般係数体 $K=k_n$、$n\ge3$ に限る。
source を $A=B/(\Omega)$ と書き、指定された右加群の自己交叉周期を用いる。
target は実際の次数三線束を持つ $Y=X_{\mathrm{res}}$、
その実際の $K$-線型 rotation を $\rho_Y$ とする。
命題\ref{prop:nfm-period-equation-cube-frame}の三乗 frame の代表を一つ選び、
それを $F_{\mathrm{sc}}$ と書く。
$\tau_Y:\rho_Y^3\simeq[2]$ はこの生成元に対する指定周期である。
以下の結果は、周期位相の値を計算する有限な入力を特定するが、
その値が一であると結論するものではない。

\begin{lemma}[自己交叉 class と三次関係の正規化]
\label{lem:nfm-tate-cubic-relation-normalization}
補題\ref{lem:nfm-cubic-tate-cut}の設定で
$E=\operatorname{Ext}_{\operatorname{Gr}A}^{*,*}(K,K)$ とし、
$t_\Omega\in E^{2,-3}$ を指定自己交叉の class とする。
$X=(E^{1,-1})^*$ と置く。
右加群の bar 規約を固定し、
Lu--Palmieri--Wu--Zhang\cite[Corollary~B]{NFMLPWZExt}の
関係と高次積の双対同定を用いる。
その符号は、中央元 $\Omega$ に対して二つの持上げ微分の合成が
$+\Omega$ となる自己交叉規約に合わせる。
このとき
\[
 R_2=\operatorname{im}\left(
       (m_2)^*:(E^{2,-2})^*\longrightarrow X^{\otimes2}\right),
 \qquad B=T_K(X)/(R_2),
\]
であり、$t_\Omega^*(t_\Omega)=1$ とすると
\[
 \Omega=\left[(m_3)^*(t_\Omega^*)\right]\quad\text{in }B_3.
\]
最後の式は、$t_\Omega$ と cohomology を固定する minimal 模型の変更に依存しない。
ここで $m_3$ を三次 tensor として gauge 不変とは主張しない。
\end{lemma}

\begin{proof}
$A$ は次数一生成であり、最小関係は二次関係 $R_2$ と、
$R_2X+XR_2$ を法とする一本の三次関係 $\Omega$ である。
同文献 Corollary~B は、最小関係の包含の双対が
$(E^{1,-1})^{\otimes d}$ に制限した $m_d$ となる模型を与える。
三次 relation line と $E^{2,-3}$ の pairing の正規化を確認すればよい。

最小 graded free resolution の最初の二つの微分を $B$ に持ち上げる。
二次関係の列では合成は零であり、三次関係 $\Omega$ の列では合成は
指定した $+\Omega$ である。
従って、この合成を $\Omega$ で割って augmentation を取る obstruction cocycle は、
三次関係の列を一、二次関係の列を零へ送る。
これは中央商の標準 truncation の connecting class の分解上の表示である。
具体的には、固定入力Bの
\texttt{lp:prop:sing-period} と \texttt{tw:eq:period-extension} が、
持上げ微分の合成を $+\Omega$ とする二拡張の Yoneda class を
まさに $t_\Omega$ と同定している。
よって $t_\Omega(\Omega)=1$ であり、上の双対公式は同じ符号で成立する。
この符号は $K[x]/(x^3)$ の
\[
 0\longrightarrow K(-3)\xrightarrow{x^2}
 A(-1)\xrightarrow{x}A\longrightarrow K\longrightarrow0
\]
でも直接確認できる。この二拡張は、最小分解の $A(-3)$ の基底を
$+1$ へ送る cocycle を表し、固定入力の MF $(x,x^2)$ の指定周期と一致する。
左右加群や bar suspension の別規約を採る場合は、
source と target の両方に同じ符号変更を施す。
片側だけに符号を入れ替えた比較は用いない。

模型の変更を、まず一次成分が恒等な $A_\infty$ 同型 $f$ で表す。
次数により $f_2(E^{1,-1},E^{1,-1})\subset E^{1,-2}=0$ である。
arity 三の同型条件において $m_3$ を変える残りの項は
\[
 f_2(m_2(-,-),-),\qquad f_2(-,m_2(-,-))
\]
だけであり、それらの双対は $R_2X+XR_2$ に入る。
従って $B_3$ における class は変わらない。
一般の一次成分は $X$ の基底変換と relation line の双対変換を与えるため、
この構成はその同型に自然である。
\end{proof}

\begin{proposition}[正準強化 class の周期位相は二次・三次積で評価できる]
\label{prop:nfm-period-phase-low-products}
$C=C_{\mathrm{res}}$、$p=p_{\mathrm{res}}$ とし、
$p$ を表す線束の Picard gerbe を
$\pi:\mathscr P_p\to\operatorname{Spec}K$、その普遍線束を $\mathcal M$ とする。
幾何的制限 $q_Y=\mathbf Li_C^*j_Y$ の右随伴を用いて
\[
 Z_p=q_Y^R\mathcal M[1]
       \in\operatorname{Ku}(Y_{\mathscr P_p})
\]
と置く。次が成り立つ。
\begin{enumerate}[label=\textup{(\roman*)}]
\item ordinary な graded orbit の全射複体
\[
 \mathcal S_p=
 \bigoplus_{q\in\mathbf Z}
   R\operatorname{Hom}(Z_p,\rho_Y^qZ_p)
\]
は、積と全 $A_\infty$ 構造を含めて $K$ 上へ降下する。
ここでは orbit の次数 $q$ を忘れない。
指定周期は、降下した自然な非零 class
\[
 t_p\in H^2(\mathcal S_p)_{-3}
\]
を与える。

\item 双次数を保つ strictly unital minimal 模型 $S_p$ を取り、
\[
 E_p^{\mathrm{ord}}=
 \bigoplus_{-i-q\ge0}S_p^{i,q},\qquad
 X_p=(S_p^{1,-1})^*,
\]
\[
 R_{p,2}=\operatorname{im}
    \bigl((m_2)^*:(S_p^{2,-2})^*\to X_p^{\otimes2}\bigr),
 \quad B_p=T_K(X_p)/(R_{p,2}),
\]
\[
 \Omega_p=\left[(m_3)^*(t_p^*)\right]\in(B_p)_3,
 \qquad t_p^*(t_p)=1,
\]
と置く。これらは minimal 模型に依存しない同型類を定める。
$(B_p,K\Omega_p)$ は $(B,K\Omega)$ の $K$-形式であり、
$\Omega_p$ は非零の中央正則三次元である。

\item $K$ 上に次数付き対の同型
\[
 f:(B_p,K\Omega_p)\xrightarrow{\sim}(B,K\Omega)
\]
が存在する。任意のそのような $f$ に対し
\begin{equation}\label{eq:nfm-period-low-product-evaluation}
 f(\Omega_p)=a_f\Omega\quad(a_f\in K^\times)
 \quad\Longrightarrow\quad
 \gamma_{\mathrm{sc},n}=[a_f]
       \in K^\times/K^{\times3}.
\end{equation}
従って位相の算出には、この対象から得る
$m_2,m_3$、指定周期 $t_p$、および次数三以下の表示の線型比較が十分である。
実際の $K$-線束 $M$ や実際の $K$-比較核の存在は仮定しない。
\end{enumerate}
\end{proposition}

\begin{proof}
\emph{orbit と周期の降下。}
$q_Y^R$ と $\rho_Y$ はいずれも $K$ 上の実際の perfect kernel 関手である。
$Z_p$ の scalar inertia は重み一なので、
$R\operatorname{Hom}(Z_p,\rho_Y^qZ_p)$ の inertia は始域と終域で相殺する。
従って各 $q$ の perfect 複体は重み零である。
$\mathbf G_m$-gerbe 上の重み零の準連接複体の fpqc 降下により、
これらは $K$ 上の複体へ降下する。
積は $\rho_Y$ の反復と射の合成で定義され、両方ともこの降下に自然である。
cofibrant な dg orbit 模型を使うか、各次数で同じ bar 模型を取れば、
全構造が降下した Adams dg algebra を得る。
これは各 Hom の cohomology の次元だけを降ろす議論ではない。

$\tau_Y$ を $\rho_Y^{-3}Z_p$ に評価した
\[
 t_p:Z_p\longrightarrow\rho_Y^{-3}Z_p[2]
\]
も同じ理由で降下する。
これは $Z_p\ne0$ 上の周期同型なので非零である。
target の生成元 $F_{\mathrm{sc}}$ とその指定周期をここで実際に用いている。
単なる $\rho_Y^3\simeq[2]$ の存在から正規化してはいない。

\emph{通常 Ext と三次関係。}
系\ref{cor:nfm-generic-adjoint-class-unique}と
定理\ref{thm:nfm-adjoint-point-finite-class-torsor}により、
この gerbe は正準強化比較 class の実現 gerbe と band を保って同値である。
その上の普遍比較核を $P$ とし、
随伴復元の同型を使えば
\[
 \Phi_P\kappa_{\mathrm{nc}}\simeq Z_p
\]
である。ここで $Z_p$ は通常の可換側の原点 stabilization ではなく、
指定した $q_Y^R\mathcal M[1]$ であることに注意する。
幾何的体拡大上で rotation の intertwiner を一つ選べば、
$\mathcal S_p$ は source の ordinary graded orbit algebra と同型となる。
補題\ref{lem:nfm-ordinary-ext-intrinsic-cut}の次数計算は、
$E_p^{\mathrm{ord}}$ が augmented $A_\infty$ 部分代数となることを与える。
同じ補題と Adams Koszul 双対により、その Koszul dual は
$(B/(\Omega))_{\overline K}$ と同型な $K$-代数となる。

補題\ref{lem:nfm-tate-cubic-relation-normalization}を適用すると、
二次関係から $B_p$、指定 $t_p$ から $B_p$ 内の三次関係 $\Omega_p$ を回復する。
模型の gauge 差が $R_{p,2}X_p+X_pR_{p,2}$ に入ることも同補題による。
幾何的体拡大後で $\Omega_p$ は $\Omega$ の非零 scalar 倍だから、
中央性、正則性、および $B_p$ の quadratic Koszul AS-regular 性は
忠実な体拡大から反映される。
特にここで使う $
 S_p^{1,-1},S_p^{2,-2},S_p^{2,-3}$ は全て有限次元で、
最後の空間は $Kt_p$ である。

\emph{元の体上の線型比較。}
次数付き対の Isom scheme は、$X_p\to B_1$ の可逆線型写像に
二次関係と三次関係直線の保存を課して得られる。
従ってこれは有限表示であり、幾何的に非空である。
補題\ref{lem:nfm-generic-presentation-stabilizer-trivial}により
\[
 \operatorname{Aut}_{\mathrm{gr}}(B,K\Omega)=\mathbf G_m,
\]
ここで右辺は次数の scalar 作用である。
従ってこの Isom scheme は $\mathbf G_m$-torsor である。
$H^1(K,\mathbf G_m)=0$ から $K$-点 $f$ を持つ。
任意の別の $f'$ は $f$ の後に次数 scalar $\lambda\in K^\times$ を合成したものであり、
\[
 a_{f'}=\lambda^3a_f.
\]
よって $[a_f]$ はこの選択にも依存しない。

\emph{位相との一致と向き。}
比較 gerbe 上の直線
$V=\operatorname{Hom}^0(P\rho_{\mathrm{nc}},\rho_YP)$ は重み零であり、
$K$ 上の直線 $V_0$ に降下する。
$0\ne e\in V_0(K)$ を取り、その引戻しを $\eta$ と書く。
これは実際の $K$-核の選択を必要としない。
普遍核と随伴復元同型を介する orbit の比較を $\Psi_\eta$ とする。
$\eta$ を正負の整数回反復することにより、
これは Adams 次数を保つ $A_\infty$ 同型となる。
全ての scalar inertia が相殺するので、この同型自身も $K$ 上に降りる。

周期位相の既存の規約は
\[
 (\tau_Y P)\eta^{(3)}=c(e)P\tau_{\mathrm{nc}}.
\]
ここで内部次数 $-3$ への orbit 比較では $\eta$ の逆反復を用いる。
従って
\begin{equation}\label{eq:nfm-period-tate-generator-comparison}
 \Psi_\eta(t_\Omega)=c(e)^{-1}t_p.
\end{equation}
例えば恒等比較に対して $\eta$ を $z$ 倍すれば、
右辺の係数は $z^{-3}$ 倍となる。
これは $c(ze)=z^3c(e)$ と一致する。

通常 Ext への制限の Koszul dual は反変であるから、
対応する同型の向きは $f_\eta:B_p\to B$ となる。
三次関係 line で \eqref{eq:nfm-period-tate-generator-comparison} を双対化すると
\[
 f_\eta(\Omega_p)=c(e)^{-1}\Omega.
\]
arity 三の同型式の補正は既述の二次関係 ideal に入るため、
この等式は $B_3$ で正確に成立する。
したがって $[a_{f_\eta}]=[c(e)^{-1}]=\gamma_{\mathrm{sc},n}$ である。
任意の $K$-同型 $f$ との差は次数 scalar だから、
\eqref{eq:nfm-period-low-product-evaluation}を得る。

最後に target だけを $vF_{\mathrm{sc}}$ へ変更すると、
$t_p$ は $v^{-1}t_p$、その双対と $\Omega_p$ は $v$ 倍となる。
従って右辺は $[v]$ 倍となり、
命題\ref{prop:nfm-period-equation-frame-change}の変換則と一致する。
source と target を同時に同じ $u$ 倍すると、分子と分母がともに $u$ 倍となる。
線型持上げを変更する target の cube 倍も Kummer 類を変えない。
\end{proof}

\begin{remark}[この評価式で残る計算と相対的な条件]
\label{rem:nfm-period-low-product-scope}
この式は、source の指定自己交叉周期、target の cube frame、
および幾何的に定義された $q_Y^R\mathcal M[1]$ を同じ有限計算に置く。
抽象的な Serre 自然性だけで trace の scalar が一になるとは主張しない。
さらに、source の stabilization を可換側の通常の原点へ送ると仮定してはいない。
その置換をすると、回復する二次包絡を誤る。

実際の計算では、$\mathscr P_p$ の有限分裂体上で $\mathcal M$ と
$q_Y^R\mathcal M[1]$ の有限 perfect 模型を作り、
内部次数 $-1,-2,-3$ の Hom、合成、および一つの contraction から
$m_2,m_3$ を求めればよい。
その結果の重み零性により、最終の二次関係束と三次 relation は元の体へ降下する。
有限行列の cohomology を求めることと、未計算の $m_3$ を零と置くことは異なる。
後者は許されない。

この一般体上の結果だけから、境界 trait 上での非分岐性は従わない。
trait 上へ用いるには、該当する Ext 空間の平坦性と底変更、
二次関係束の一定階数、$t_p$ が規格化する三次 relation の原始性、
および表示の射影 Isom の延長を別に確認する必要がある。
また、orbit の低次積の回復は coherent な中心線型構造の降下そのものではない。
\end{remark}

\begin{lemma}[良い integral 表示による周期位相の非分岐判定]
\label{lem:nfm-period-integral-presentation-unramified}
$R$ を $\mathbb C$-代数である離散付値環、$K=\operatorname{Frac}(R)$ とする。
source の次数付き対 $(B,R\Omega)$ と、
命題\ref{prop:nfm-period-phase-low-products}の有限低積から回復した対
$(B_p,R\Omega_p)$ が、次の条件を満たすとする。
\begin{enumerate}
\item 両者の次数一生成元加群 $V,V_p$ は同じ階数の有限自由 $R$-加群であり、
両対は第\ref{sec:poisson-inverse}節の良い表示開 $J^{\mathrm{sm}}$ の $R$-対象である。
すなわち、$T^{\mathrm{alg}}$ に含めた相対 PBW・普遍中心・相対 Serre 双対・
全 $\mathcal O$-同値比較の入力、境界評価 $\epsilon$ の同型条件、
相対点スキームの良い条件、および trace--Hessian 三次の滑らかさを、
\emph{閉ファイバーを含めて}満たす。
ここで $\Omega,\Omega_p$ は、それぞれ次数三部分の直線部分束の生成元である。
\item 有限周期評価に現れる射影化された次数付き同型
\[
 \bar f_K:(B_p,K\Omega_p)\xrightarrow{\sim}(B,K\Omega)
       \quad\bmod\mathbf G_m
\]
は、幾何点だけでなく $K$ 上の点として与えられている。
その任意の線型持上げ $f_K$ について、式\eqref{eq:nfm-period-low-product-evaluation}により、考える位相の評価式が
\[
 \gamma=[u_K],\qquad
 f_K(\Omega_p)=u_K\Omega
       \quad\text{in }B_3\otimes_RK
\]
である。式の三次関係は二次関係による剰余で読む。
\end{enumerate}
このとき $\gamma$ は $R$ 上の有限 \'etale な $\mu_3$-torsor へ延びる。
すなわち
\[
 v_R(\gamma)=0\quad\text{in }\mathbb Z/3.
\]
この判定は、比較核そのものの $R$ 上への延長を仮定しない。
\end{lemma}

\begin{proof}
まず、ここで使用する有限性の対象を特定する。
定理\ref{thm:poisson-good-open-isomorphism}は、上の二つの良い対を
$\mathfrak P_n^{\mathrm{good}}\subset\mathfrak P_n^{\mathrm{dec}}$ の二つの対象へ移し、
次数スカラーで割った次数付き同型を、その full Poisson 対の同型と同定する。
滑らかな三次と超平面偏極類の射影同型空間は有限であり、
Poisson tensor を保つ条件はその中の閉条件である。
従って補題\ref{lem:nfm-poisson-target-finite-diagonal}の証明は、
符号商を取らないこの同型空間にも適用される。
特に
\[
 I=\underline{\operatorname{Isom}}_R
       \bigl((B_p,R\Omega_p),(B,R\Omega)\bigr)/\mathbf G_m
\]
は $R$ 上有限である。これは線型同型の空間が有限という主張ではない。

指定された $K$-点 $\bar f_K$ は有限射 $I\to\operatorname{Spec}R$ の
有理的切断である。$R$ は正規なので、この切断は $R$ 上へ一意に延長する。
例えば $I=\operatorname{Spec}A$ と書けば、$A\to K$ の像は $R$ 上整であり、
$R$ に含まれる。得られた切断を $\bar f_R$ とする。

両次数一加群を自明化し、$\bar f_R$ の ambient 射影同型を
$\operatorname{PGL}_N(R)$ の点と見る。
fppf 完全列
\[
 1\longrightarrow\mathbf G_m\longrightarrow\operatorname{GL}_N
       \longrightarrow\operatorname{PGL}_N\longrightarrow1
\]
から、その線型持上げの空間は $R$ 上の $\mathbf G_m$-torsor であり、
障害は $H^1(R,\mathbf G_m)=\operatorname{Pic}(R)=0$ に入る。
従って $f_R\in\operatorname{Isom}_R(V_p,V)$ を選べる。
これは一般の Brauer--Severi 空間の分裂を $\operatorname{Pic}(R)=0$ から
推論しているのではない。最初から二つの有限自由加群と、その射影空間間の
\emph{実際の}射影同型が与えられている場合の持上げである。
次数スカラーは二次関係と中心直線を保つから、$f_R$ は両対の次数付き同型になる。

特に $f_R(R\Omega_p)=R\Omega$ であり、両者の指定生成元について
\[
 f_R(\Omega_p)=u_R\Omega,\qquad u_R\in R^\times
\]
を得る。仮定で指定した $f_K$ と $f_R|_K$ は同じ射影同型の持上げなので、
ある $\lambda\in K^\times$ を用いて $f_K=\lambda f_R|_K$ と書ける。
次数三上では $u_K=\lambda^3u_R$ であり、$[u_K]=[u_R]$ である。
従って $\operatorname{Spec}R[z]/(z^3-u_R)$ が所要の有限 \'etale な延長を与える。
\end{proof}

\begin{remark}[この判定で残る閉ファイバー条件]
\label{rem:nfm-period-integral-presentation-scope}
この補題は、$\Delta D\ne0$ だけからその仮定を導くものではない。
有限 perfect な orbit Hom 模型が延びることと、その低 cohomology が
$R$-自由で底変更と可換になることは異なる。
さらに、通常 Ext の切出しに必要な積と homotopy transfer を integral に行えること、
二次関係束が期待階数を持つこと、指定周期から回復した三次関係
$\Omega_p$ が primitive であること、および回復表示が上の
$J^{\mathrm{sm}}$ に残ることを別に確認する必要がある。
有限 perfect spread は、これらが成立する非空開を与えても、
指定した余次元一ファイバーをその開が横切るとは限らない。
また、$\bar f_K$ の幾何的存在だけでは上の有理切断を得ない。
有限拡大 $K'/K$ の後でのみ適用する場合には、その DVR 上の非分岐性から戻す際に
$v_{R'}(a)=e(R'/R)v_R(a)$ を用い、分岐指数が三で割れる場合は元の付値を判定できない。


現行の保証開 $T$ の構成では、active 点移動に
$\sigma^j\ne1$、$1\le j\le\max(3,n-1)$ を課している。
従って $n\ge5$ では、$\Delta D=0$ とは別に現れる
$\sigma$ が位数四となる divisor なども、現在の入力の保証範囲から外れる。
この active 条件は混合係数を一般にしても変わらない。
そこでの PBW・普遍中心・生成性が実際に破れると断定するものではないが、
現在の入力をそのまま用いて、全ての $\Delta D$ 外の divisor で
本補題の閉ファイバー条件が成立すると結論することはできない。
\end{remark}
